\chapter{ชุดพรอมพต์วิศวกรรม 5 ระดับสำหรับการพัฒนา}
\label{chap:ch04}

\section*{แผนบริหารการสอนประจำบทที่ 4}
\addcontentsline{toc}{section}{แผนบริหารการสอนประจำบทที่ 4}

\begin{tcolorbox}[
    enhanced, breakable,
    colback=white, colframe=rbruNavy,
    fonttitle=\bfseries\color{white},
    title={1. หัวข้อเนื้อหาประจำบท},
    arc=1.5mm, boxrule=0.8pt, leftrule=3.5mm
]
\begin{enumerate}
    \item กรอบการออกแบบชุดคำสั่งตามหลักการ CLEAR Prompting Framework
    \item การออกแบบสถาปัตยกรรมระดับระบบและโครงสร้างโปรเจกต์ (System Prompting)
    \item การสร้างตัวถอดรหัสโปรโตคอล Modbus RTU สำหรับอ่านค่า pH อุณหภูมิ และความชื้น
    \item การเปลี่ยนแบบจำลองคณิตศาสตร์ฟิสิกส์สู่โค้ดโครงข่ายประสาทเทียม PINN บน Dart
    \item เทคนิคการสร้างหน้าจอ UI แบบ Responsive ป้องกันปัญหาข้อมูลล้นจอทุกมิติ
\end{enumerate}
\end{tcolorbox}

\begin{tcolorbox}[
    enhanced, breakable,
    colback=white, colframe=rbruNavy,
    fonttitle=\bfseries\color{white},
    title={2. วัตถุประสงค์เชิงพฤติกรรม},
    arc=1.5mm, boxrule=0.8pt, leftrule=3.5mm
]
เมื่อศึกษาบทเรียนนี้จบแล้ว ผู้เรียนสามารถ
\begin{enumerate}
    \item อธิบายองค์ประกอบ 5 ประการของ CLEAR Prompting ได้อย่างถูกต้อง
    \item ออกแบบชุดคำสั่งสำหรับสั่งการโมเดลปัญญาประดิษฐ์เพื่อสร้างโค้ดแอปพลิเคชัน SoilpH-HT ได้อย่างแม่นยำ
    \item ประยุกต์ใช้เทคนิค Few-Shot Prompting ในการเขียนชุดทดสอบ Unit Test ได้อย่างสมบูรณ์
    \item วินิจฉัยและแก้ไขข้อผิดพลาดในการรันแอปพลิเคชันผ่านกระบวนการ Prompt Chaining ได้อย่างเป็นระบบ
\end{enumerate}
\end{tcolorbox}

\begin{tcolorbox}[
    enhanced, breakable,
    colback=white, colframe=rbruNavy,
    fonttitle=\bfseries\color{white},
    title={3. วิธีสอนและกิจกรรมการเรียนการสอน},
    arc=1.5mm, boxrule=0.8pt, leftrule=3.5mm
]
\begin{enumerate}
    \item การสาธิตการป้อนพรอมพต์วิศวกรรมและการวิเคราะห์ผลตอบกลับของโมเดล AI
    \item กิจกรรมฝึกเขียนพรอมพต์เพื่อแก้ปัญหาโจทย์การตรวจวัดเซนเซอร์จริงแบบรายบุคคล
    \item การทดลองรันโค้ดที่สร้างขึ้นบนโปรแกรมจำลองและแก้ไขข้อผิดพลาดร่วมกันในชั้นเรียน
\end{enumerate}
\end{tcolorbox}

\begin{tcolorbox}[
    enhanced, breakable,
    colback=white, colframe=rbruNavy,
    fonttitle=\bfseries\color{white},
    title={4. สื่อการเรียนการสอน},
    arc=1.5mm, boxrule=0.8pt, leftrule=3.5mm
]
\begin{enumerate}
    \item เอกสารคู่มือชุดคำสั่งพรอมพต์วิศวกรรมและเทมเพลตมาตรฐาน
    \item ระบบผู้ช่วยปัญญาประดิษฐ์สำหรับการเขียนโปรแกรม (AI Coding Assistant)
    \item สภาพแวดล้อมการพัฒนา Flutter SDK และ Visual Studio Code
\end{enumerate}
\end{tcolorbox}

\begin{tcolorbox}[
    enhanced, breakable,
    colback=white, colframe=rbruNavy,
    fonttitle=\bfseries\color{white},
    title={5. การวัดและประเมินผล},
    arc=1.5mm, boxrule=0.8pt, leftrule=3.5mm
]
\begin{enumerate}
    \item การประเมินโครงสร้างและความสมบูรณ์ของชุดคำสั่งพรอมพต์ที่ผู้เรียนออกแบบ
    \item การประเมินผลการทำงานจริงของโค้ดโปรแกรมที่สร้างจากชุดคำสั่ง
    \item การทดสอบความเข้าใจเกี่ยวกับกลวิธี CLEAR Prompting
\end{enumerate}
\end{tcolorbox}

\clearpage

% ==============================================================================
% ภาพประกอบเกริ่นนำประจำบทที่ 4 (Introductory Infographic)
% ==============================================================================
\begin{figure}[H]
\centering
\begin{tikzpicture}[
    box/.style={rectangle, rounded corners=2mm, draw=rbruNavy, line width=1.0pt, fill=white, drop shadow={shadow xshift=1pt, shadow yshift=-1.5pt, opacity=0.15}, align=center, inner sep=5pt},
    header/.style={rectangle, rounded corners=1.5mm, draw=rbruNavy, fill=rbruNavy, text=white, font=\bfseries\scriptsize, align=center, inner sep=3.5pt},
    arrow/.style={-{Stealth[scale=1.0]}, line width=1.1pt, draw=rbruBlue}
]
    % วงจร CLEAR 5 ขั้นตอน
    \node[box, fill=cyan!12, minimum width=2.4cm] (c) {\textbf{C - Context}\\[2pt]\scriptsize กำหนดบทบาท\\และบริบทระบบ};
    \node[box, fill=rbruGold!15, minimum width=2.4cm, right=0.5cm of c] (l) {\textbf{L - Logic}\\[2pt]\scriptsize ฝังกฎฟิสิกส์\\และสมการหลัก};
    \node[box, fill=green!12, minimum width=2.4cm, right=0.5cm of l] (e) {\textbf{E - Example}\\[2pt]\scriptsize ตัวอย่างเฟรม\\และข้อมูลเข้าออก};
    \node[box, fill=orange!15, minimum width=2.4cm, right=0.5cm of e] (a) {\textbf{A - Action}\\[2pt]\scriptsize สั่งการสร้างโค้ด\\ตามลำดับขั้น};
    \node[box, fill=purple!12, minimum width=2.4cm, right=0.5cm of a] (r) {\textbf{R - Result}\\[2pt]\scriptsize โค้ด Dart/Flutter\\พร้อมชุดทดสอบ};

    \draw[arrow] (c) -- (l);
    \draw[arrow] (l) -- (e);
    \draw[arrow] (e) -- (a);
    \draw[arrow] (a) -- (r);

    % ลำดับขั้น 5 ระดับด้านล่าง
    \node[rectangle, rounded corners=2mm, draw=rbruNavy!70, fill=rbruNavy!6, line width=0.8pt, minimum width=14.5cm, minimum height=1.0cm, below=0.8cm of e, align=center] (ladder) {
        \small\textbf{ลำดับขั้นการพัฒนา 5 ระดับ (5-Level Progressive Engineering Pipeline)}\\[2pt]
        \footnotesize
        1. สถาปัตยกรรมระบบ $\rightarrow$ 2. เฟรม Modbus \& CRC-16 $\rightarrow$ 3. โมเดล PINN สองขั้นตอน $\rightarrow$ 4. มาตรวัด pH แบบ Responsive $\rightarrow$ 5. ไดรเวอร์ USB-C OTG
    };
    \draw[arrow, dashed] (e.south) -- (ladder.north);
\end{tikzpicture}
\caption{ผังกรอบแนวคิดวิศวกรรมพรอมพต์แบบ CLEAR และบันไดการพัฒนา 5 ระดับสำหรับการสร้างแอปพลิเคชัน SoilpH-HT}
\label{fig:ch04_prompt_pipeline}
\end{figure}

\section{กรอบการออกแบบชุดคำสั่งตามหลักการ CLEAR Prompting}

การใช้ปัญญาประดิษฐ์เชิงกำเนิดในการพัฒนาซอฟต์แวร์ระดับระบบและการเชื่อมต่อฮาร์ดแวร์จริง จำเป็นต้องอาศัยชุดคำสั่งที่มีความรัดกุม แม่นยำ และมีบริบทที่ชัดเจน กรอบการทำงาน \textbf{CLEAR Prompting Framework} ประกอบด้วย 5 มิติหลัก ได้แก่
\begin{enumerate}
    \item \textbf{Context (บริบท)} ระบุบทบาท หน้าที่ และสภาพแวดล้อมทางเทคโนโลยี เช่น วิศวกรสมองกลฝังตัวและนักวิทยาศาสตร์ข้อมูลฟิสิกส์เกษตร
    \item \textbf{Logic (ตรรกะและหลักการ)} กำหนดสมการ กฎทางฟิสิกส์ หรือมาตรฐานที่ต้องปฏิบัติตาม
    \item \textbf{Example (ตัวอย่างประกอบ)} ให้ตัวอย่างรูปแบบข้อมูลเข้า (Input) และข้อมูลออก (Output) ที่ต้องการ
    \item \textbf{Action (การกระทำที่ชัดเจน)} สั่งการสิ่งที่ต้องการให้ทำเป็นข้อๆ อย่างชัดเจน
    \item \textbf{Result (รูปแบบผลลัพธ์)} กำหนดโครงสร้างผลลัพธ์ เช่น ซอร์สโค้ดภาษา Dart พร้อมชุดทดสอบ
\end{enumerate}

\section{พรอมพต์ระดับที่ 1 การออกแบบสถาปัตยกรรมและโครงสร้างโฟลเดอร์}

\begin{promptbox}
\begin{lstlisting}
คุณคือสถาปนิกซอฟต์แวร์ระบบสมองกลฝังตัวและการเกษตรแม่นยำ (AgriTech Software Architect)
จงออกแบบโครงสร้างโปรเจกต์แอปพลิเคชันมือถือ Flutter สำหรับแอปพลิเคชัน SoilpH-HT ตรวจวัดและวิเคราะห์ค่ากรด-ด่างดิน
ชดเชยอุณหภูมิและความชื้นด้วย Deep Learning PINN ผ่านพอร์ต USB Type-C OTG
โดยมีข้อกำหนดดังต่อไปนี้
1. จัดโครงสร้างแบบ Clean Architecture แบ่งเป็น core, data, domain, ui (MVVM)
2. กำหนดเนื้อหาไฟล์ pubspec.yaml พร้อม dependencies ที่จำเป็น
3. กำหนดเนื้อหาไฟล์ AndroidManifest.xml พร้อมสิทธิ์ android.hardware.usb.host และการเข้าถึงกล้องและ GPS
4. กำหนดไฟล์ device_filter.xml รองรับ Vendor ID ของชิปแปลง USB ยอดนิยม ได้แก่ CH340, CP2102, FTDI, PL2303
5. ออกแบบให้มีโหมดจำลองข้อมูล (Demo Simulator) เมื่อยังไม่ได้เสียบอุปกรณ์จริง
\end{lstlisting}
\end{promptbox}

\section{พรอมพต์ระดับที่ 2 การเขียนคลาสถอดรหัส Modbus RTU}

\begin{promptbox}
\begin{lstlisting}
คุณคือวิศวกรระบบสื่อสารอุตสาหกรรม (Industrial Protocols Engineer)
จงเขียนคลาสภาษา Dart ชื่อ ModbusParser สำหรับถอดรหัสข้อมูลจากหัววัด Soil Parameter Tester มีข้อกำหนดดังนี้
1. ฟังก์ชัน calculateCRC16 คำนวณแบบบิตต่อบิตตามมาตรฐาน Modbus RTU (Polynomial 0xA001, ค่าเริ่มต้น 0xFFFF)
2. ฟังก์ชัน buildReadCommand สร้างชุดคำสั่ง 8 ไบต์ ร้องขออ่าน Holding Register ของความชื้น อุณหภูมิ EC และ pH
3. ฟังก์ชัน parseResponse ตรวจสอบความยาวเฟรม ตรวจสอบ CRC-16 และแปลงข้อมูล 16-bit Big-Endian เป็นตัวเลขจริง
4. เขียน Unit Test ตรวจสอบความถูกต้องของฟังก์ชันคำนวณ CRC-16
\end{lstlisting}
\end{promptbox}

\section{พรอมพต์ระดับที่ 3 การพัฒนาโมเดล Deep Learning PINN}

\begin{promptbox}
\begin{lstlisting}
คุณคือนักวิทยาศาสตร์ข้อมูลฟิสิกส์เกษตรดิจิทัล (Agriphysics AI Data Scientist)
จงพัฒนาคลาสภาษา Dart ชื่อ DeepLearningCalibrator สำหรับรันโมเดล Physics-Informed Neural Network ชดเชยค่า pH ดิน
โดยมีเงื่อนไขสำคัญดังนี้
1. ไม่พึ่งพาไลบรารีภายนอก (Pure Dart Forward-Pass Inference) เพื่อให้ทำงานได้เร็วกว่า 60 FPS
2. นำทฤษฎีฟิสิกส์มาใช้ชดเชย ได้แก่ Nernstian Slope Drift, Water Dissociation Shift (pKw) และ Dry-Soil Contact Impedance
3. ฝังโครงข่ายประสาทเทียม MLP 3 ชั้น ใช้ฟังก์ชันกระตุ้น Swish และ GELU ชดเชยปฏิสัมพันธ์ร่วมระหว่างอุณหภูมิและความชื้น
4. คืนค่าเป็น Map ประกอบด้วยค่าดิบ (Raw), ค่าที่ชดเชยแล้ว (Calibrated) และค่าผลต่าง (Delta)
\end{lstlisting}
\end{promptbox}

\section{พรอมพต์ระดับที่ 4 การพัฒนา UI มาตรวัด pH แบบ Responsive}

\begin{promptbox}
\begin{lstlisting}
คุณคือนักออกแบบ Flutter UI/UX ขั้นสูง (Senior Flutter UI Engineer)
จงออกแบบหน้าจอแดชบอร์ดหลักของแอปพลิเคชัน SoilpH-HT มีข้อกำหนดดังนี้
1. จัดวางมาตรวัดวงกลม PhCircularGauge เป็นจุดเด่นตรงกลาง ปรับขนาดตามหน้าจอด้วย LayoutBuilder อัตโนมัติ (170 ถึง 240 px)
2. ป้องกันปัญหา Layout Overflow อย่างเด็ดขาดบนสมาร์ตโฟนจอแคบ (< 360 px) ด้วย FittedBox, Expanded และ ConstrainedBox
3. มีแผงเทเลเมทรีแสดงอุณหภูมิ ความชื้น และค่าผลต่างการชดเชย Delta pH
4. มีปุ่มเปิดกล้องถ่ายภาพดินพร้อมเป้าเล็งและประทับพิกัดดาวเทียม GPS และข้อมูลเทเลเมทรีลงบนภาพถ่าย
5. เมื่อแตะที่มาตรวัด ให้เปิด Bottom Sheet แสดงผลการวินิจฉัยดินและคำแนะนำอัตราปูนโดโลไมต์ปรับปรุงดิน
\end{lstlisting}
\end{promptbox}

\section{พรอมพต์ระดับที่ 5 การเชื่อมต่อพอร์ต USB Serial บนแอนดรอยด์}

\begin{promptbox}
\begin{lstlisting}
คุณคือนักพัฒนาระบบฝังตัวบนแอนดรอยด์ (Android Low-Level Embedded Developer)
จงเขียนคลาส UsbPhService ภาษา Dart เพื่อจัดการการเชื่อมต่อ USB OTG Serial กับหัววัด Soil Parameter Tester
1. ตรวจจับการเสียบและถอดสาย USB OTG อัตโนมัติ (Hotplug Detection)
2. ขออนุญาตสิทธิ์การเข้าถึงอุปกรณ์ผ่านระบบ Android USB Manager
3. มีระบบสลับอัตราเร็ว Baud rate อัตโนมัติ (4800 หรือ 9600 bps)
4. มีระบบเฝ้าระวัง Watchdog Timer ขนาด 3 วินาที กู้คืนการเชื่อมต่ออัตโนมัติเมื่อหัววัดไม่ตอบสนอง
5. สลับเข้าสู่โหมดจำลอง (Simulator Mode) อัตโนมัติเมื่อเกิดข้อผิดพลาดในการเชื่อมต่อ
\end{lstlisting}
\end{promptbox}

\section{สรุปสาระสำคัญประจำบทที่ 4}

วิศวกรรมพรอมพต์ (Prompt Engineering) มิใช่เพียงการพิมพ์คำถามทั่วไป แต่คือการกำหนดกรอบความคิดทางวิศวกรรมซอฟต์แวร์ให้แก่แบบจำลองภาษาขนาดใหญ่ โดยสรุปสาระสำคัญได้ดังนี้
\begin{enumerate}
    \item \textbf{กรอบการทำงาน CLEAR} การระบุ Context, Logic, Example, Action และ Result อย่างครบถ้วน ช่วยตัดปัญหาการสร้างโค้ดที่ผิดพลาด (Hallucination) และทำให้ AI สร้างโค้ดที่ถูกต้องตามหลักฟิสิกส์และมาตรฐานความปลอดภัย
    \item \textbf{การพัฒนาแบบก้าวหน้า 5 ระดับ} การแบ่งการสั่งการออกเป็น 5 ชั้นโมดูลาร์ (สถาปัตยกรรม $\rightarrow$ การสื่อสารข้อมูล $\rightarrow$ โมเดลคณิตศาสตร์ $\rightarrow$ ส่วนต่อประสานกราฟิก $\rightarrow$ ไดรเวอร์ฮาร์ดแวร์) ช่วยให้ควบคุมคุณภาพของซอฟต์แวร์ได้อย่างรัดกุม
    \item \textbf{การทดสอบกำกับในทุกขั้นตอน} ทุกพรอมพต์ต้องระบุให้ AI สร้างชุดทดสอบหน่วย (Unit Tests) ประกอบมาด้วยเสมอ เพื่อให้สามารถตรวจสอบความถูกต้องของฟังก์ชันได้ทันที
\end{enumerate}

ในบทถัดไป (บทที่ 5) จะนำเสนอซอร์สโค้ดฉบับสมบูรณ์ที่สร้างขึ้นจากชุดพรอมพต์ทั้ง 5 ระดับ พร้อมคำอธิบายการทำงานในแต่ละโมดูลอย่างละเอียด

\section*{คำถามทบทวนท้ายบทที่ 4}
\addcontentsline{toc}{section}{คำถามทบทวนท้ายบทที่ 4}

\begin{enumerate}
    \item จงอธิบายบทบาทของขั้นตอน Context และ Logic ในกรอบการทำงาน CLEAR Prompting
    \item หากต้องการให้ปัญญาประดิษฐ์สร้างฟังก์ชันการถอดรหัสสัญญาณที่มีโครงสร้างข้อมูลซับซ้อน ควรใช้เทคนิค Few-Shot Prompting อย่างไร
    \item จงวิเคราะห์สาเหตุที่การสั่งการโมเดล AI โดยไม่มีการระบุกรอบสถาปัตยกรรม (Architecture Constraints) มักนำไปสู่ปัญหาโค้ดที่มีความซ้ำซ้อนและบำรุงรักษายาก
    \item จงออกแบบชุดคำสั่งสำหรับสั่งการโมเดล AI เพื่อให้ช่วยสร้างฟังก์ชันคำนวณการปรับปรุงความเป็นกรด-ด่างของดินด้วยปูนโดโลไมต์
\end{enumerate}
